\documentclass{article}
\usepackage[pdftex]{graphicx}
\usepackage{xcolor}
\usepackage[letterpaper, margin = 2.5cm]{geometry}
\usepackage[T2A]{fontenc}
\usepackage[english, russian]{babel}
\usepackage{amsmath, amsfonts, amssymb}
\usepackage{hyperref, url}

\usepackage{minted}

\bibliography{refs.bib}

\usepackage{parskip}

\begin{document}

\title{Домашнее задание 3}
\date{}
\author{SMR-2026 Starkit Education, Илья Осокин @elijahmipt} %поменяйте на своё
\maketitle

\section*{Задача 1}

\noindent 

Рассмотрим систему Cart-Pole (маятника на каретке), где масса каретки составляет $M$, масса груза маятника $m$, длина маятника $l$, и каретка двигается горизонтально в поле тяготения $g$.

\begin{enumerate}
    \item Какой у этой системы вектор состояния?
    \item Найдите лагранжиан системы
    \item С помощью уравнения Эйлера-Лагранжа найдите уравнения движения системы
    \item Линеаризуйте их в окрестности неустойчивого положения равновесия маятника
\end{enumerate}

В классрум нужно будет загрузить один .ipynb файл, проще всего будет дописать 2 и 3 задачи в конец семинарского ноутбука и загрузить их одним файлом. Кто не доделал предыдущее задание, доделайте и загрузите в classroom.

\section*{Задача 1}

Проработайте весь семинарский ноутбук, выполните задания, впишите код везде, где нужно.

\section*{Задача 2}

1) Реализуйте стабилизацию системы, состоящей из двух обычных (не обратных) маятников и пружины между ними, с помощью LQR. Основные шаги тут следующие:
\begin{itemize}
    \item Линеаризуйте систему в окрестности положения равновесия (это уже есть)
    \item Определите весовые матрицы Q и R правильных размерностей
    \item Решите алгебраическое уравнение Риккати с помощью scipy
    \item Встройте обратную связь в цикл симуляции системы. Постройте графики углов, скоростей и момента
    \item Исследуйте влияние весов в матрицах Q и R на траектории системы. Если управление дорогое, система стабилизируется быстро или медленно? Что изменится, если домножить обе матрицы Q и R на 10?
\end{itemize}

\section*{Задача 3}

Сделайте то же самое, что и в задании 2, только в окрестности верхнего положения равновесия. Попробуйте поменять весовые матрицы. Меняется ли принципиально поведение системы? Чем объясняются отличия траекторий этой системы и системы из двух маятников, висящих вниз?

Сделайте то же самое, но для линеаризованного Cart-pole.

\section*{Задача 4}

Доразберитесь со всем, что было непонятно из первых трех пар. Составьте список вопросов по теории управления, которые хочется обсудить, и напишите его в чат. Вопросы для самодиагностики (писать ответы не надо, просто подумайте):
\begin{itemize}
    \item Приведите пример нелинейной системы, которую сложно линеаризовать
    \item На какой частоте должна работать система управления автоматической машины? Автономного квадрокоптера?
    \item Какие бывают типы ограничений в задачах управления? Например, могут быть ограничены моменты, могут быть углы, а еще что?
\end{itemize}

\end{document}
